%% ma: L10-B14-0005
%% lop: 10
%% chuong: 5
%% bai: 14
%% dang: 10.14.3
%% loai: TL
%% muc: VD
%% dap_an: "-"
%% nguon: "chuyen-de-cac-so-dac-trung-cua-mau-so-lieu-khong-ghep-nhom-toan-10-chuong-trinh-moi.pdf (GV. Phan Nhật Linh) (P04), Bài 03 Dạng 1 Bài tập tự luận 5"
%% kien_thuc: "Lớp 10 Bài 13 (tứ phân vị), Bài 14 (khoảng tứ phân vị, giá trị bất thường)"
%% trang_thai: chinh-thuc
%% ghi_chu: "Tự duyệt theo yêu cầu giáo viên 22/10/2026."
\begin{ex}
Điều tra một số học sinh về số cái bánh chưng mà gia đình mỗi bạn sử dụng trong dịp Tết Nguyên đán, kết quả ghi lại ở bảng sau:
\begin{center}
\begin{tabular}{|c|c|c|c|c|c|c|c|}\hline
Số bánh chưng & $6$ & $7$ & $8$ & $9$ & $10$ & $11$ & $15$\\\hline
Số gia đình & $5$ & $7$ & $10$ & $8$ & $5$ & $4$ & $1$\\\hline
\end{tabular}
\end{center}
Mẫu số liệu trên có giá trị bất thường không?
\loigiai{$n=40$; tần số tích lũy $5$; $12$; $22$; $30$; $35$; $39$; $40$.

$Q_2=\dfrac{x_{20}+x_{21}}2=8$; $Q_1=\dfrac{x_{10}+x_{11}}2=7$; $Q_3=\dfrac{x_{30}+x_{31}}2=\dfrac{9+10}2=9{,}5$.

$\Delta_Q=2{,}5$; $Q_3+1{,}5\Delta_Q=13{,}25$; $Q_1-1{,}5\Delta_Q=3{,}25$.

Giá trị $15>13{,}25$ nên $15$ là giá trị bất thường.}
\end{ex}
